\section{Conclusion}

In this work, we presented \textbf{AniUnFlow-SamT}, an unsupervised segment-guided multi-frame optical flow framework tailored to the distinctive motion and appearance characteristics of 2D animation. In contrast to conventional two-frame unsupervised pipelines, the proposed approach integrates temporal slot memory, segment-guided object correspondence, layered affine motion priors, and coarse-to-fine dense correlation refinement with boundary-aware residual correction, thereby coupling object-level structural reasoning with dense motion recovery in a unified architecture. Trained without ground-truth flow supervision through a composite objective that combines photometric reconstruction, smoothness regularization, forward--backward consistency, segment-level coherence, boundary-aware constraints, and cycle consistency, AniUnFlow-SamT is designed to produce sharper motion boundaries, stronger intra-object coherence, and improved temporal stability in challenging anime-style scenarios characterized by texture sparsity, large non-rigid motion, and frequent occlusion changes. Overall, these results support the effectiveness of combining object-centric priors with dense correspondence learning for animation-oriented optical flow, and suggest promising directions for future research in longer-horizon temporal modeling, robustness to extreme stylization, and tighter integration with downstream animation tasks such as frame interpolation, line-art inbetweening, and temporally consistent stylized video generation.